Interpretation of the Slower Oscillation as a Higgs Mode: I question whether the observed slower oscillation can be strictly attributed to a Higgs mode (i.e., an amplitude fluctuation). If the ordered cycloid component were fluctuating, one would expect the ferroelectric polarization to both rise and fall. However, the data shows that the SHG signal repeatedly decreases and recovers, but does not increase beyond its baseline. The authors should address why the expected "rise" in polarization is absent if this is indeed an amplitude fluctuation mode. 
